\begin{table}[!htbp]\centering\small
\begin{tabular}{@{}lrrrr@{}}\toprule
Familia & núcleo tras reducir & comp.\ a PD & comp.\ a B\&R & nodos B\&R (media) \\\midrule
Erdős–Rényi $n{=}200$ & 23.0\% & 0 & 5 & 538 \\
Erdős–Rényi $n{=}500$ & 23.3\% & 0 & 2 & 391 \\
$k$-árbol parcial $n{=}300$ & 0.0\% & 0 & 0 & 0 \\
Banda $n{=}400$ & 93.7\% & 4 & 3 & 59 \\
Barabási–Albert $n{=}1000$ & 33.3\% & 0 & 2 & 4771 \\
Trampa de Johnson & 0.0\% & 0 & 0 & 0 \\
Erdős–Rényi denso $n{=}60$ & 100.0\% & 0 & 3 & 150 \\
\bottomrule\end{tabular}
\caption{Qué sección de \hyvc{} hace el trabajo: fracción de vértices que sobrevive a las reducciones, componentes despachadas a programación dinámica o a branch \& reduce, y nodos de búsqueda.}\label{tab:hyvc}\end{table}
